% ===================================================================
% Title page
% ===================================================================
\begin{titlepage}
\centering
\vspace*{2cm}

{\color{ruleblue}\rule{0.6\textwidth}{1.2pt}}\par
\vspace{0.6cm}

{\normalfont\Huge\bfseries Lean for Working Algebraists\par}
\vspace{0.5cm}
{\normalfont\Large\itshape A Guide to Formalizing Groups, Rings, and\\[2pt]
Path Algebras in Lean~4\par}

\vspace{0.6cm}
{\color{ruleblue}\rule{0.6\textwidth}{1.2pt}}\par

\vspace{2.5cm}
{\normalfont\Large Abderrahim Adrabi\par}
\vspace{0.2cm}
{\normalfont\normalsize\itshape Faculty of Sciences, Rabat, Morocco\\
Faculty of Educational Sciences, Rabat, Morocco\par}

\vfill

{\normalfont\large\itshape "Every proof in this book was verified against
a real Lean~4 toolchain --- \\
not just described."\par}

\vspace{1.5cm}
{\normalfont\normalsize \today\ --- \bookversion\par}
{\normalfont\small Lean toolchain \texttt{v4.33.1}\par}
\vspace{0.3cm}
{\normalfont\small Source code:
\url{https://github.com/abderrahim-lectures/lean4-learning}\par}

\end{titlepage}

% ===================================================================
% Copyright / license page (verso of the title page)
% ===================================================================
\clearpage
\thispagestyle{empty}
\null\vfill
\begin{center}
{\normalfont\small
Copyright \textcopyright\ 2026 Abderrahim Adrabi\par
\vspace{0.4cm}
Licensed under the MIT License. See the \texttt{LICENSE} file in the
source repository for the full license text.\par
\vspace{0.4cm}
Source code, issue tracker, and the latest release:\\
\url{https://github.com/abderrahim-lectures/lean4-learning}\par
}
\end{center}
\vfill\null
\clearpage

% ===================================================================
% Preface
% ===================================================================
\chapter*{Preface}

This is a guide to Lean~4 for readers who already think in groups, rings,
functors, and diagrams, but have never written a line of Lean --- or any
proof assistant. No programming background is assumed, only the
mathematical maturity of someone comfortable with abstract algebra and the
basic language of category theory: objects, morphisms, composition,
functors. The book builds up Lean~4 syntax and tactics only as far as
needed to formalize groups, rings, and quiver path algebras; the
categorical viewpoint --- a quiver's path category, a ring as a one-object
preadditive category, and so on --- is called out explicitly wherever it
clarifies what the Lean code is really encoding.

This book is, and remains, Mathlib-free by design through Chapter~5: every
group, ring, and path algebra is built from scratch, so every definition
and proof obligation stays visible. Starting in Chapter~6, each worked
example is followed by a ``Mathlib equivalent'' box showing the same
statement phrased against Mathlib's real API --- not a contradiction of
the from-scratch approach, but a second, parallel track. Seeing the same
idea twice, once built by hand and once as a library already has it, is
how both halves of working in Lean get learned at once: the mathematics
itself, and the shape of the library a reader will actually use
afterward. Chapter~13 completes the handoff to Mathlib in full.

The running goal throughout is more than the constructs covered: it is
the \emph{skill} of using Lean --- reading a goal state, deciding what to
try next, recovering when a tactic fails, and knowing which proofs to
derive by hand versus hand off to automation.

% ===================================================================
% About this book
% ===================================================================
\chapter*{About this book}
This is an introduction to Lean~4 for mathematicians, built from scratch.
It presupposes no programming background and no previous contact with any
proof assistant; it presupposes only the mathematical maturity of someone
comfortable with abstract algebra and the basic language of category
theory. The book builds up Lean~4 syntax and tactics only as far as needed
to formalize groups, rings, modules, and quiver path algebras from
scratch. The categorical viewpoint --- the path category of a quiver as the
free category on a quiver, a ring as a one-object preadditive category ---
is made explicit wherever it clarifies what the Lean code is really
encoding, rather than bolted on as decorative commentary.

Every concept and every theorem in the book is derived in place, and every
code block is verified against a real Lean~4 toolchain via a companion
project that compiles it with \texttt{lake build}. The mathematics is
presented as a process of derivation and discovery, not as a finished
skeleton to memorize; see the Pedagogical approach section that follows.

\chapter*{Pedagogical approach}
This book teaches by derivation, not by dispensation, in the tradition of
the teaching of Arnold and Gelfand rather than that of Bourbaki.
``It is impossible to understand an unmotivated definition'' (Arnold,
\emph{On Teaching Mathematics}, 1997). Every definition and theorem in
this book is \emph{earned}. A section poses the question or concrete
problem that forces a concept, walks the reasoning that discovers it, and
only then names and formalizes the result, mirroring the method of the
Moscow seminar Gelfand ran, which finds the simplest example that
captures a phenomenon before finding the language to state it generally.
No fixed
Definition--Theorem--Proof--Remark skeleton is imposed on that reasoning;
structure follows the logic of the argument, not a template. Formulas and
theorems are proved before they are used, never handed down as rules to
memorize. Exercises favor fewer, harder, escalating, proof-heavy problems
over repetitive drills, and do not give away their own answer; full
solutions live in the appendix.

% ===================================================================
% A note on authorship
% ===================================================================
\chapter*{A note on authorship}
This book was written as an experiment in using Claude Code, the AI
coding agent built by Anthropic, to produce a complete technical book.
The author, Abderrahim Adrabi --- a mathematician (PhD) and
techno-pedagogical engineer affiliated with the Faculty of Sciences,
Rabat, and the Faculty of Educational Sciences, Rabat, Morocco ---
directed the scope, structure, and audience, and has reviewed the
mathematical and Lean content directly, chapter by chapter, across every
round of revision; Claude wrote the chapter text, Lean code, and LaTeX.
No additional, independent mathematician or Lean expert has reviewed the
content beyond direct review by the author and the automated checks
described in the README and CHANGELOG of the project (every code block
compiled against a real Lean~4 toolchain, plus several AI-driven review
passes).

Each release has gone through further closing verification passes,
described in full in \texttt{reviews/} and in the project CHANGELOG.
These passes re-triage prior findings, check citations against the
literature, and re-derive worked examples and proofs by hand rather than
accepting them as written. Confirmed issues are corrected before release
and recorded in the changelog; unconfirmed predictions made by a
reviewer are never applied without independent verification.

Treat this book as an author-reviewed draft produced through
human-directed AI authorship, not a traditionally peer-reviewed text.

% ===================================================================
% How to read this book
% ===================================================================
\chapter*{How to read this book}

Not every reader needs every chapter in order. The Learning paths chapter
just after this one lays out a chapter-dependency graph and a handful of
named paths through the book --- already know Lean? already know algebra?
want the formal foundations first? --- for readers who would rather not
start cover to cover.

Each chapter is divided into short, numbered sections building on one
another. Every code listing is valid Lean~4 (toolchain
\texttt{v4.33.1}). Every code block in Chapters~1--11 has additionally
been ported into a companion Lean project and verified with
\texttt{lake build} against the real compiler, not merely read over ---
several bugs that only surfaced this way (a missing extensionality
lemma, a couple of \texttt{rw} steps that over-rewrote, a reference to an
undefined name, a tactic used without its library ever being imported)
have been fixed in both the book and the project.

This book is about more than the constructs it covers. Chapters~7 and~9
present each theorem as a search process --- what to look at, what to
try, why an attempt fails --- rather than only the polished final proof;
Chapter~12 is dedicated entirely to working efficiently once the
underlying ideas are understood. Chapter~5 pauses to address the rigor
questions a careful mathematician will already be asking by that point
--- \texttt{structure} versus \texttt{class}, universes, and definitional
versus propositional equality --- before Chapter~6 commits to
\texttt{Group}'s definition.

A notation reference follows this section, for quick lookup as unfamiliar
symbols are encountered; two further lookup tables --- a tactic and
library reference, and a $\lambda$-calculus-to-Lean dictionary --- sit at
the back of the book, after the last chapter, to be consulted as needed
rather than read start to finish.

\tableofcontents
